\chapter{Algebraic complexity classes} \label{chap-vpvnp}

Valiant \cite{v79} defined two algebraic complexity classes that can be thought of as analogues of the boolean classes $\P$ and $\NP$. 
This chapter focuses on their definitions, and some important properties related to the polynomial families $\Det_n$ and $\Perm_n$. 

\section{Definition of the Basic Classes}

Recall that $\P$ is\footnote{If one is to be more precise, this is $\P/\poly$ or non-uniform $\P$. 
But in this article, we shall be interested only in the non-uniform versions since we mainly deal with circuit sizes.} the class of boolean functions that can be computed by circuits of polynomial size. 
As any boolean function can be expressed as a multilinear\footnote{A polynomial where the degree in any variable is bounded by $1$.} polynomial, an analogue of this in the arithmetic world could be multilinear polynomials $f(x_1,\dots, x_n)$ that can be computed by arithmetic circuits of size $\poly(n)$. 
However, unlike in the boolean world, the polynomial $x^2$ is not equal to the polynomial $x$ as we are dealing with formal polynomials. 
Valiant's definition of $\VP$ was the class of the class of ``low degree'' polynomials that can be computed by circuits of small size. 

\begin{definition}[Valiant's $\P$]\label{defn:vp}
The class $\VP$ refers to the set of polynomials $f(x_1,\dots, x_n)$ of \emph{degree $\poly(n)$} that can be computed by arithmetic circuits of size $\poly(n)$. 
\end{definition}

In the literature, one also encounters classes such as $\mathsf{VBP}$ and $\mathsf{VF}$ that correspond to polynomials computed by polynomial-sized ABP and formulas respectively. 
These are subclasses of $\VP$ by definition. 


\subsection*{More on the low-degree restriction}

But should the analogue of $\VP$ not be the class of polynomials that are computable by $\poly(n)$-sized arithmetic
circuits, including polynomials of very large degree?
We can indeed compute polynomials of very large degree, such as a circuit that is a chain of $s$ multiplication gates thus computing a polynomial of degree $\mathrm{exp}(s)$ (by repeated squaring). 
Let us first take a moment to understand why the additional restriction of ``low-degree'' in the above definition of $\VP$ was imposed. 
There are several intuitive reasons for this, and this ``restricted'' definition also yields beautiful structural results. 
This discussion is from an answer by Joshua Grochow on \texttt{cstheory.SE} \cite{gro:SE} to shed more light on the above definition. 

\begin{enumerate}

\item Every boolean function can be expressed as a multilinear polynomial. 
Multilinear polynomials are, of course, polynomials of ``low degree''. 

\item Much of the earlier work was based on understanding formula size. 
In the case of arithmetic formulas, the degree cannot be more than the size of the formula. 
If a polynomial is computed by a $\poly(n)$ sized formula, then its degree must be bounded by $\poly(n)$ too. 

\item Most interesting polynomials, such as $\Det_n$ or $\Perm_n$, are in fact of low degree. 
Once we choose to deal with only polynomials of low degree, the above definition does not have any restriction on the circuit used to compute it (besides its size). 
It is certainly possible that intermediate computations of the circuit could involve very large degree polynomials. 

However, as we shall soon see, Strassen's result shows that such high degree computations may be eliminated with just an $O(\deg^2)$ increase in size. 
Dealing with low-degree circuits also makes the class robust under this transformation. 
Eliminating division gates also only incurs an $O(\deg^2)$ increase in size. 

\item Without this restriction, one cannot hope to show that $\Det_n$ or $\Perm_n$ is ``complete'' for such classes, or
show the numerous structural results such as depth reduction that we have. 

\end{enumerate}

Having said that, there is also a notion of ``degree'' in a boolean circuit that is defined syntactically as follows:
 \begin{quote}
   Degree of all leaves is $1$. 

   For any OR gate, the degree is the maximum of the degree of its children. 

   For any AND gate, the degree is the sum of the degree of its children. 
 \end{quote}
The class of boolean functions that can be computed by poly-sized poly-degree circuits coincide with a class called $\mathsf{LOGCFL}$ or $\mathsf{SAC}^1$. 
With this notion, one might say that $\VP$ is really an analogue of $\mathsf{SAC}^1$.\\


We now move on to the arithmetic analogue of the class $\NP$. 
Recall that the class $\NP$ may be defined as the set of all boolean functions $f(x_1,\dots, x_n)$ such that there is some $g(x_1,\dots, x_n, y_1,\dots, y_m)$ with $m = \poly(n)$ and 
\[
f(x_1,\dots, x_n)\spaced{=} \bigvee_{\veca \in \inbrace{0,1}^m}  g(x_1,\dots, x_n, a_1,\dots, a_m).
\]
Valiant's $\NP$ is defined analogously by replacing the OR by a sum. 

\begin{definition}[Valiant's $\NP$]\label{defn:vnp}
The class $\VNP$ is defined to be the set of polynomials $f(x_1,\dots, x_n)$ such that there is some $g(x_1,\dots, g_n, y_1,\dots, y_m) \in \VP$ with $m = \poly(n)$ and
\[
f(x_1,\dots, x_n)\spaced{=} \sum_{\veca \in \inbrace{0,1}^m}  g(x_1,\dots, x_n, a_1,\dots, a_m).
\]
\end{definition}

It follows from definitions that $\mathsf{VF} \subseteq \mathsf{VBP} \subseteq \VP \subseteq \VNP$. 
We do not know if any of the containments is strict (although it is widely believed that all of them are). 

%\section{Definition of Border Classes (?)}

\section{Some properties of these classes}

We shall state a few properties of these classes here. 
For a more extensive treatment, B\"{u}rgisser's book \cite{bur00} has a comprehensive study of these classes and a lot of the proofs presented in this chapter are based on the description in his book. \\

Valiant~\cite{v79} presented a very useful sufficient condition to show that a given polynomial is in $\VNP$. 

\begin{theorem}[Valiant's Criterion] \label{thm:val-criterion}
Let $f = \sum_\vece c(\vece) \cdot x_1^{e_1}\dots x_n^{e_n}$. 
Suppose the function $\varphi$ that takes as input the exponent vector $\vece = (e_1,\dots, e_n)$ and outputs the coefficient $c(\vece)$ is in the class $\#\P/\poly$, then $f \in \VNP$. 
\end{theorem}

Thus in particular, if we compute the coefficient for a given monomial in polynomial time, then the polynomial is in $\VNP$. 

\begin{corollary}
The polynomials $\Perm_n$ and $\NW_{n,m,d}$ are in $\VNP$. 
\end{corollary}

As stated in \autoref{defn:vnp}, the polynomial $g(x_1,\dots, x_n, y_1,\dots, y_m) \in \VP$. 
However, a subsequent paper of Valiant~\cite{v82} showed that we may assume without loss of generality that $g \in \mathsf{VF}$, that is $g$ is computable by a small formula. 
This is similar to the fact that counting solutions of 3-\textsf{CNF} instance, which is a formula, is as hard as counting solutions of any polynomial sized boolean circuit. 
We state this result here, and a proof may be found in B\"urgisser's book \cite{bur00}. Malod and Portier \cite{MP08} presented an alternate proof via \emph{proof-trees} which is perhaps more instructive. 

\begin{theorem}\label{thm:vnp-formula}
For any $f(x_1,\dots, x_n)\in \VNP$, there is a $g(x_1,\dots, x_n, y_1,\dots, y_m)$ that can be computed by a $\poly(n)$ sized \emph{formula} such that 
\[
f(x_1,\dots, x_n)\spaced{=} \sum_{\veca \in \inbrace{0,1}^m} g(x_1,\dots, x_n, a_1,\dots, a_m).
\]
\end{theorem}

\subsection{Closure under factoring (?)}

Just mention the results without proof? 
